\medskip
\hypertarget{native-109-article080}{}\pdfbookmark[2]{80.}{article-080}
\begin{center}80.\end{center}
If the angle $\Pi$ is given, $p$ and $e$ are determined by the equations
\[
\begin{aligned}
p&=\frac{rr'(\cos(N-\Pi)-\cos(N'-\Pi))}{r\cos(N-\Pi)-r'\cos(N'-\Pi)}\\[4pt]
e&=\frac{r'-r}{r\cos(N-\Pi)-r'\cos(N'-\Pi)}\mathpunct{\cdot}
\end{aligned}
\]
The common denominator in these formulas can be put in the form $a\cos(A-\Pi)$ with $a$ and $A$ independent of $\Pi$. Indeed, if $H$ denotes an arbitrary angle, we have
\[
\begin{aligned}
&r\cos(N-\Pi)-r'\cos(N'-\Pi)\\
&\quad=(r\cos(N-H)-r'\cos(N'-H))\cos(H-\Pi)\\
&\qquad-(r\sin(N-H)-r'\sin(N'-H))\sin(H-\Pi),
\end{aligned}
\]
and thus the denominator is $=a\cos(A-\Pi)$ if $a$ and $A$ are determined by the equations
\SourceLeaf{110}
\[
\begin{aligned}
r\cos(N-H)-r'\cos(N'-H)&=a\cos(A-H)\\
r\sin(N-H)-r'\sin(N'-H)&=a\sin(A-H).
\end{aligned}
\]
This gives
\[
\begin{aligned}
p&=\frac{2rr'\sin\tfrac12(N'-N)\sin(\tfrac12N+\tfrac12N'-\Pi)}{a\cos(A-\Pi)}\\[4pt]
e&=\frac{r'-r}{a\cos(A-\Pi)}\mathpunct{\cdot}
\end{aligned}
\]
These formulas are particularly convenient whenever $p$ and $e$ are to be calculated for several values of $\Pi$ while $r$, $r'$, $N$, and $N'$ are held fixed. — For calculating the auxiliary quantities $a$ and $A$, the angle $H$ may be chosen arbitrarily. It is convenient to set $H=\tfrac12(N+N')$. The equations for these quantities then become
\[
\begin{aligned}
(r'-r)\cos\tfrac12(N'-N)&=-a\cos(A-\tfrac12N-\tfrac12N')\\
(r'+r)\sin\tfrac12(N'-N)&=-a\sin(A-\tfrac12N-\tfrac12N').
\end{aligned}
\]
Having thus determined the angle $A$ from the equation
\[
\operatorname{tang}(A-\tfrac12N-\tfrac12N')=\frac{r'+r}{r'-r}\operatorname{tang}\tfrac12(N'-N),
\]
we immediately have
\[
e=-\frac{\cos(A-\tfrac12N-\tfrac12N')}{\cos\tfrac12(N'-N)\cos(A-\Pi)}\mathpunct{\cdot}
\]
The calculation of the logarithm of $\dfrac{r'+r}{r'-r}$ can be shortened by the device already explained several times.
